% ID: FIS-ELE-POT-001
% Titolo: Potenziale elettrico generato da una carica puntiforme
% Disciplina: Fisica
% Area: Elettrostatica
% Argomento: Potenziale elettrico
% Sottoargomento: Differenza di potenziale e lavoro del campo
% Classe: Quinta liceo scientifico
% Difficolta: 3
% Tipo: problema_numerico
% Risultato: V_A = 60 V; V_B = 30 V; L = 1,5 * 10^{-7} J
% Tempo_stimato: 9 min
% Competenze: applicazione della formula del potenziale, uso della differenza di potenziale, interpretazione energetica
% Prerequisiti: carica elettrica, campo elettrico, lavoro di una forza
% Tag: fisica, elettrostatica, potenziale_elettrico, carica_puntiforme, lavoro
% Fonte: esempio_originale
% Autore: OpenAI Codex
% Licenza: CC-BY-SA-4.0

\begin{esercizio}
Una carica puntiforme positiva \(Q=2{,}0\,\mathrm{nC}\) genera un campo
elettrico nello spazio circostante. Considera due punti \(A\) e \(B\) posti
rispettivamente a \(30\,\mathrm{cm}\) e \(60\,\mathrm{cm}\) dalla carica.

Calcola il potenziale elettrico nei due punti e il lavoro compiuto dal campo
elettrico per spostare una carica positiva \(q=5{,}0\,\mathrm{nC}\) da
\(A\) a \(B\). Usa \(k=9{,}0\cdot 10^9\,\mathrm{N\,m^2/C^2}\).
\end{esercizio}

\begin{soluzione}
Il potenziale generato da una carica puntiforme e:
\[
V=\frac{kQ}{r}.
\]

Per il punto \(A\):
\[
V_A=\frac{9{,}0\cdot 10^9\cdot 2{,}0\cdot 10^{-9}}{0{,}30}
=60\,\mathrm{V}.
\]

Per il punto \(B\):
\[
V_B=\frac{9{,}0\cdot 10^9\cdot 2{,}0\cdot 10^{-9}}{0{,}60}
=30\,\mathrm{V}.
\]

Il lavoro compiuto dal campo elettrico nello spostamento da \(A\) a \(B\) e:
\[
L=q(V_A-V_B).
\]

Sostituendo:
\[
L=5{,}0\cdot 10^{-9}\cdot (60-30)
=1{,}5\cdot 10^{-7}\,\mathrm{J}.
\]

Il lavoro e positivo perche la carica positiva si sposta da un potenziale
maggiore a un potenziale minore, cioe nel verso favorito dal campo.
\end{soluzione}
